\chapter{Nouns and the Noun Phrase}\label{ch:nouns}

\section{Introduction}

Nominal morphology in Iridian is relatively straightforward. Nouns are inflected for case; there is no grammatical gender, and number marking is optional. The case system interacts closely with the verbal aspect system: the morphological role of a noun phrase in a clause depends in part on the aspect of the verb, a consequence of Iridian's split-ergative alignment (discussed in detail in Chapter~\ref{ch:verbs}).

\section{The case system}

Iridian has three productive cases: the \term{direct}, the \term{transitive}, and the \term{indirect}. A fourth case, the \term{vocative}, exists for a restricted set of nouns used in direct address. The direct case is the default, uninflected form from which all other cases are derived.

The functions of each case are closely tied to the aspect of the verb in the clause:

\begin{itemize}
	\item The \term{direct case} (\Dir{}) is the unmarked or citation form of a noun. It functions as the unmarked or core argument of the clause: the subject of an intransitive verb, the subject of a transitive verb in the nominative-alignment aspects (continuous, progressive, contemplative), and the object of a transitive verb in the transitive-alignment aspects (perfective and retrospective). It is also used as the predicative complement in equative constructions (zero copula) and with the overt copular verb \ird{neh-}; and as the entity whose existence is asserted in existential constructions with \ird{jec-} and \ird{zad-}.

	\item The \term{transitive case} (\Trs{}) is the ergative case. It marks the agent of a transitive verb in the perfective and retrospective aspects. Outside of those contexts, it does not appear as a core argument.

	\item The \term{indirect case} (\Ind{}) is the oblique or non-core case. It marks the patient of a transitive verb in the nominative aspects (continuous, progressive, contemplative); it also marks instruments, beneficiaries, locations, and other peripheral arguments regardless of aspect. Its function is thus both grammatical and adverbial.

	\item The \term{vocative} (\Voc{}) is used for direct address. It exists only for personal names and a small set of common nouns, primarily kinship terms.
\end{itemize}

\pex
\a
\begingl
\gla Marek zmacime.//
\glb Marek.\Dir{} run-\Prog{}//
\glft \trsl{Marek is running.} \lingcontext{intransitive; agent = \Dir{}}//
\endgl
\a
\begingl
\gla Travat Marku piaštnik.//
\glb bread.\Dir{}.\Def{} Marek.\Trs{} eat-\Pf{}//
\glft \trsl{Marek ate the bread.} \lingcontext{transitive perfective; P = \Dir{}, A = \Trs{}}//
\endgl
\a
\begingl
\gla Marek travě piastime.//
\glb Marek.\Dir{} bread.\Ind{} eat-\Prog{}//
\glft \trsl{Marek is eating bread.} \lingcontext{transitive progressive; A = \Dir{}, P = \Ind{}}//
\endgl
\xe

\section{Inflectional classes}

Nouns in Iridian are best described through two intersecting systems: \term{declension class} and \term{stem grade}. The declension class determines which case endings a noun selects; the stem grade determines which shape of the stem appears before that ending. This allows the grammar to treat forms such as \ird{Marek} $\rightarrow$ \ird{Marku} / \ird{Marcí}, \ird{obel} $\rightarrow$ \ird{oblu} / \ird{oblí}, and \ird{lobek} $\rightarrow$ \ird{lubku} not as isolated curiosities, but as regular consequences of the general morphophonemic system described in Chapter~\ref{ch:phon}.

The surface case inventory remains simple --- direct, transitive, indirect, and vocative --- but the stems that feed those cases may occur in several recurrent grades:

\begin{itemize}
	\item the \term{basic stem}, used in the direct/citation form;
	\item the \term{zero-grade stem}, where weak \ird{e} is deleted before an oblique ending;
	\item the \term{front-raised stem}, where zero-grade formation further yields \ird{e/ě/é $\rightarrow$ i/í};
	\item the \term{back-raised stem}, where zero-grade formation yields \ird{o $\rightarrow$ u};
	\item the \term{soft stem}, where a front-vocalic ending palatalizes the final consonant;
	\item the \term{leveled stem}, where an inherited or high-frequency form overrides the productive pattern.
\end{itemize}

Not every noun shows all these grades, but every morphophonemically complex noun can be described by reference to them. Dictionaries should therefore identify nouns not merely by a citation form, but by a compact inflectional profile: \term{Direct form}, \term{declension class}, \term{weak-\ird{e} status}, \term{raising behavior}, and any \term{leveled indirect or vocative form}. In ordinary prose discussion this grammar cites the direct form alone unless some further principal part is needed.

Native direct-case nouns may end in a consonant, in \ird{-a}, in \ird{-o}, in \ird{-e/-ě}, or in \ird{-ou}. Final \ird{-i} and \ird{-u} are avoided in inherited absolutive nouns and occur chiefly in recent loans, abbreviations, or marginal expressive vocabulary. The declension system must therefore account not only for consonant-final stems, but also for several distinct vowel-final root shapes.

Modern standard Iridian has eighteen declension classes. These are lexically assigned and must be learned noun by noun, but they are not arbitrary: they fall into six broader macrofamilies, each with three internal subclasses. Within a given class the endings and stem alternations are regular; what is lexical is chiefly which class a noun belongs to, whether it tolerates deletion and raising, and whether it preserves an inherited short indirect ending or one of the longer areally remodeled endings.

\begin{table}[ht]
	\scriptsize\sffamily
	\caption{The eighteen noun declension classes.}\label{tab:declension}
	\medskip
	\begin{tabularx}{\textwidth}{l L l l l L}
		\toprule
		{\sc class} & {\sc type} & {\sc \Trs{}} & {\sc \Ind{}} & {\sc \Voc{}} & {\sc stem behavior} \\
		\midrule
		1 & hard consonant plain & \ird{-u} & \ird{-í} & \ird{-o} & basic stem; no obligatory weak-\ird{e} deletion \\
		2 & hard consonant deleting & \ird{-u} & \ird{-ení} & \ird{-o} & productive weak-\ird{e} deletion plus expanded oblique ending \\
		3 & hard consonant front-raised & \ird{-u} & \ird{-í} & \ird{-o} & deletion plus \ird{e/ě/é $\rightarrow$ i/í} \\
		4 & hard consonant back-raised & \ird{-u} & \ird{-ení} & \ird{-e} & deletion plus \ird{o $\rightarrow$ u}; long front-vocalic indirect \\
		5 & hard animate arealized & \ird{-a} & \ird{-ovi} & \ird{-e} & Slavic-like animate oblique endings beside inherited stem grades \\
		6 & hard consonant leveled & \ird{-u} & lexical \ird{-í / -ení} & zero/\ird{-o} & productive stem behavior, but ending choice partly analogical \\
		7 & velar plain & \ird{-u} & \ird{-e} & \ird{-e} & regular velar affrication in \Ind{} and \Voc{} \\
		8 & velar deleting & \ird{-u} & \ird{-ení} & \ird{-e} & weak-\ird{e} deletion plus velar softening under an expanded indirect ending \\
		9 & velar arealized & \ird{-a} & \ird{-evi} & \ird{-e} & animate-like oblique remodelled under areal pressure \\
		10 & soft consonant plain & \ird{-y} & \ird{-ě} & \ird{-e} & already soft stem; little further alternation \\
		11 & soft consonant contracted & \ird{-y} & \ird{-i} & \ird{-e} & contracted front oblique with little further alternation \\
		12 & sibilant / affricate expanded & \ird{-y} & \ird{-ení} & \ird{-e} & class-specific spellout after sibilants and affricates \\
		13 & \ird{-a} stem & \ird{-y} & \ird{-ě} & --- & low alternation; default native vocalic class \\
		14 & \ird{-o} stem (\ird{o + m > um}) & \ird{-a} & \ird{-um} & --- & final \ird{-o} replaced in \Trs{}; coalesces with nasal indirect marker \\
		15 & front-vowel stem in \ird{-e/-ě} & \ird{-e} & \ird{-ení} & --- & levelled front-vowel stems with analogically extended indirect \\
		16 & consonant \ird{-a/-am} & \ird{-a} & \ird{-am} & \ird{-e} & conservative oblique family with little softening \\
		17 & weak-\ird{e} \ird{-a/-am} & \ird{-a} & \ird{-am} & \ird{-e} & productive deletion before oblique endings, often with syllabic liquids \\
		18 & contracting \ird{-ou} & \ird{-ovy} & \ird{-ově / -ovení} & --- & inherited contraction plus occasional areal extension of the indirect \\
		\bottomrule
	\end{tabularx}
\end{table}

The classes are numerous because Iridian noun morphology is intentionally noun-heavy. The same four cases recur everywhere, but the path from citation form to oblique form is richly differentiated by class. Some classes preserve short inherited endings, others generalize longer front-vocalic endings such as \ird{-ení}, and a few animate classes have borrowed or convergently renewed endings such as \ird{-ovi} and \ird{-evi}. The language thus preserves an agglutinative syntax and transparent case inventory while allowing a much denser Slavic-looking nominal surface.

\subsection{Macrofamily I: hard consonant stems}

Classes~1--6 comprise the hard consonant family. Here the microfamilies show genuine ending divergence as well as stem divergence. Class~1 preserves the short inherited indirect \ird{-í}. Class~2 extends that oblique to \ird{-ení}, giving forms of the type \ird{louž} $\rightarrow$ \ird{loužení}. Classes~3 and~4 add front-vowel and back-vowel raising respectively, but differ again in whether the indirect is the short \ird{-í} or the expanded \ird{-ení}. Class~5 contains animate nouns remodeled under areal pressure, with \ird{-a} in the transitive and \ird{-ovi} in the indirect. Class~6 keeps a partly leveled oblique paradigm in high-frequency nouns and personal names.

Representative outputs are:

\pex
\a \ird{Marek} $\rightarrow$ \ird{Marku} / \ird{Marcí} \lingcontext{short inherited indirect}
\a \ird{louž} $\rightarrow$ \ird{loužu} / \ird{loužení} \lingcontext{expanded indirect in Class~2}
\a \ird{lobek} $\rightarrow$ \ird{lubku} / \ird{lubcení} \lingcontext{back-raised stem plus long indirect}
\a \ird{Marek} in the arealized animate class $\rightarrow$ \ird{Marka} / \ird{Markovi}
\xe

\subsection{Macrofamily II: velar stems}

Classes~7--9 are a specialized subset of consonant-final nouns whose final consonant is velar. They preserve the visible velar alternations that most strongly index the language's Slavic surface profile. In this family the key analysis is velar affrication: after weak-\ird{e} deletion, /k/ surfaces as \ird{c} before front-vocalic oblique endings. Class~7 takes a short indirect in \ird{-e}; Class~8 uses the longer \ird{-ení}; Class~9 has animate-like areal endings in \ird{-evi}. The distinction among the three classes thus lies not only in deletion and raising, but also in which oblique ending the affricated velar stem selects.

\pex
\a \ird{perek} $\rightarrow$ \ird{prku} / \ird{prce}
\a \ird{človek} $\rightarrow$ \ird{človku} / \ird{človcení}
\a \ird{rok} $\rightarrow$ \ird{roka} / \ird{rocevi}
\xe

\subsection{Macrofamily III: soft consonant stems}

Classes~10--12 contain stems that already end in a soft consonant, a sibilant, or an affricate. Class~10 preserves the inherited \ird{-ě} indirect. Class~11 contracts that ending to \ird{-i}. Class~12 generalizes an expanded \ird{-ení}, often under areal pressure and especially after sibilants. This is the chief source of contrasts such as short \ird{Marcí} beside longer \ird{loužení}.

Because these nouns are already soft or sibilant-final, they usually show less dramatic alternation than hard consonant nouns. Their role in the system is not to showcase weak-\ird{e} deletion, but to preserve a dense inherited declensional texture.

\subsection{Macrofamily IV: vocalic stems}

Classes~13--15 are the main vowel-final families. Class~13 contains the default \ird{-a} stems such as \ird{trava}, with transitive \ird{-y} and indirect \ird{-ě}. Class~14 contains back-vowel \ird{-o} stems, whose transitive replaces the final vowel with \ird{-a} and whose indirect shows the regular coalescence \ird{o + m > um}; thus \ird{bro} $\rightarrow$ \ird{bra} / \ird{brum}. Class~15 contains front-vowel stems in \ird{-e/-ě}; these are often historically complex and are more heavily leveled than the other classes, and many of them now show an analogically extended indirect in \ird{-ení}. Together with the separate \ird{-ou} class below, these vocalic families account for the full set of native vowel-final absolutive shapes.

\pex
\a \ird{trava} $\rightarrow$ \ird{travy} / \ird{travě}
\a \ird{bro} $\rightarrow$ \ird{bra} / \ird{brum}
\a \ird{země} $\rightarrow$ \ird{zeme} / \ird{zemení}
\xe

\subsection{Macrofamily V: the \ird{-a/-am} oblique family}

Classes~16 and~17 preserve the old consonant-final family whose transitive is \ird{-a} and indirect \ird{-am}. Class~16 is conservative and comparatively stable. Class~17, by contrast, subjects the stem to productive weak-\ird{e} deletion before those oblique endings and often resolves the resulting cluster through a syllabic liquid rather than by restoring the vowel. The especially productive subtype is the \term{CeLeC} pattern, where the direct stem contains weak \ird{e} before a liquid and a final consonant, and the oblique stem surfaces as \ird{CLC-} plus the case vowel licensed by its class.

\pex
\a \ird{byl} $\rightarrow$ \ird{byla} / \ird{bylam}
\a \ird{jelen} $\rightarrow$ \ird{jlna} / \ird{jlnam}
\a \ird{berel} in the \ird{CeLeC} subtype $\rightarrow$ \ird{brla} / \ird{brlam}
\a \ird{berel} in the most reduced pronunciation $\rightarrow$ [br̩lam]
\xe

\subsection{Macrofamily VI: the contracting \ird{-ou} class}

Class~18 contains inherited \ird{-ou} nouns, above all deverbal nominals. These follow a distinct but internally regular pattern: transitive \ird{-ovy}, indirect \ird{-ově}, and contraction in the direct definite form. Some members, especially literary or formal nouns, also admit an analogically extended indirect in \ird{-ovení}. This class is highly visible in nominalizations and therefore serves as the chief bridge between noun-heavy morphophonology and the rest of the grammar.

\pex
\a \ird{piaščikou} $\rightarrow$ \ird{piaščikovy} / \ird{piaščikově}
\a \ird{hladou} $\rightarrow$ \ird{hladovy} / \ird{hladovení}
\a \ird{piaščikou + -ot} $\rightarrow$ \ird{piaščikovat}
\xe

\subsection{Case endings and stem selection}\label{sec:noun-stem-grades}

The relation between case and stem grade is summarized below.

\begin{table}[ht]
	\footnotesize\sffamily
	\caption{Case endings and their morphophonemic effects.}\label{tab:noun-stem-selection}
	\medskip
	\begin{tabularx}{0.95\textwidth}{L L L}
		\toprule
		{\sc case slot} & {\sc typical endings} & {\sc stem selected} \\
		\midrule
		Direct & zero & basic stem \\
		Transitive & \ird{-u}, \ird{-y}, \ird{-e}, \ird{-a}, \ird{-ovy} & basic, zero-grade, raised, or vowel-replacing stem according to declension class \\
		Indirect & \ird{-í}, \ird{-e}, \ird{-ení}, \ird{-ovi}, \ird{-evi}, \ird{-ě}, \ird{-i}, \ird{-um}, \ird{-am}, \ird{-ově}, \ird{-ovení} & soft stem for front-vocalic endings; zero-grade or raised stem where the class licenses deletion; some classes replace final root vowels or add areally expanded oblique material \\
		Vocative & \ird{-o}, \ird{-e}, zero & basic or soft stem; often levelled in names and kin terms \\
		\bottomrule
	\end{tabularx}
\end{table}

The transitive and indirect cases are therefore not merely suffixes added to an invariant base. Rather, the ending first selects a stem grade, and only then surfaces in its overt form. The rule order is the same as in Chapter~\ref{ch:phon}: suffix class selection first, then weak-\ird{e} deletion or retention, then any raising, then soft-stem selection if the ending is front-vocalic, and only after that the cluster and voicing repairs that yield the final phonological form.

Irregularity in the noun system is best understood as \emph{class assignment plus stem-grade irregularity}, not as unconstrained suppletion. Some nouns fail to raise where the general pattern would suggest it, some preserve weak \ird{e} in only one oblique, and some classes have generalized longer endings such as \ird{-ení} or animate-like endings such as \ird{-ovi} under areal pressure. Once these lexical facts are known, however, the resulting paradigms are still learnable by class rather than item by item.

\pex
\a \ird{Marek} $\rightarrow$ \ird{Marku} / \ird{Marcí}
\a \ird{perek} $\rightarrow$ \ird{prku} / \ird{prce}
\a \ird{bro} $\rightarrow$ \ird{bra} / \ird{brum}
\a \ird{louž} $\rightarrow$ \ird{loužu} / \ird{loužení}
\a \ird{obel} $\rightarrow$ \ird{oblu} / \ird{oblí}
\a \ird{lobek} $\rightarrow$ \ird{lubku}
\a \ird{piaščikou} $\rightarrow$ \ird{piaščikovy} / \ird{piaščikově}
\xe

\subsection{Irregular nouns}

A small number of nouns have irregular inflectional patterns that do not conform neatly even to the eighteen-class system. These are best treated as explicitly levelled lexemes. Dictionaries accordingly list any exceptional indirect or vocative form when it cannot be recovered from the noun's class and grade behavior alone.

\section{The vocative}

The vocative is used for direct address. It is productive for personal names and for a closed set of common nouns, mostly kinship terms. Outside this set, nouns are addressed with the direct case form.

Nouns with vocative forms include the following kinship terms:

\begin{table}[ht]
	\footnotesize\sffamily
	\caption{Kinship terms with vocative forms.}\label{tab:vocative}
	\medskip
	\begin{tabularx}{0.7\textwidth}{L l l}
		\toprule
		{\sc direct} & {\sc vocative} & {\sc meaning} \\
		\midrule
		\ird{papka}  & \ird{papko}  & father \\
		\ird{mamka}  & \ird{mamko}  & mother \\
		\ird{mlazka} & \ird{mlazko} & brother \\
		\ird{voska}  & \ird{vosko}  & sister \\
		\bottomrule
	\end{tabularx}
\end{table}

For personal names the vocative is formed in a manner predictable from the inflectional class of the name (see Table~\ref{tab:declension} for the hard consonant plain type exemplified by \ird{Marek}, vocative \ird{Marko}). The vocative list is not exhaustive; other common nouns may have vocatives, but their use has become increasingly restricted to formal and literary registers.

\section{Definiteness}\label{sec:definiteness}

Iridian possesses a postposed determiner, which appears as \ird{-ot} after a consonant and as \ird{-t} after a vowel. Thus \ird{bylot} signifies \trsl{the child}, and \ird{travat} \trsl{the bread}. This suffix is not to be understood as answering in every particular to the English definite article. Its office is rather to mark a noun as referring to something already present to the mind of the speaker and hearer, whether by previous mention, by the circumstances of the utterance, or by some other sort of immediate intelligibility. It is therefore more exact to say that the suffix marks an \textit{anchored} or \textit{determinate} reference than that it merely marks definiteness in the narrower sense.

In form the determiner is attached to the fully inflected noun. Since the direct case is without overt ending, the suffix there appears immediately upon the nominal stem; where, however, the noun bears some case termination, the determiner follows that termination. One may accordingly set side by side such forms as \ird{byl} \trsl{child}, \ird{bylot} \trsl{the child}, \ird{bylam} \trsl{child} in the indirect case, and \ird{bylamot} \trsl{the child} in the same case.

\ex
\begingl
\gla byl \qquad bylot \qquad bylam \qquad bylamot//
\glb child.\Dir{} \qquad child.\Dir{}.\Def{} \qquad child.\Ind{} \qquad child.\Ind{}.\Def{}//
\glft \trsl{child, the child, child (oblique), the child (oblique)}//
\endgl
\xe

The use of the determiner is optional. The bare noun is the ordinary and least marked form, and may be employed where the reference is indefinite, generic, kind-denoting, weak, or otherwise not yet fixed to the common ground of discourse. The suffixed form, on the other hand, presents the referent as one which is in some manner already recoverable. Very often it is anaphoric; not seldom it is situationally unique; at other times it is contrastively selected or resumed as a discourse-prominent participant. The opposition is therefore not simply that between \trsl{definite} and \trsl{indefinite}, but rather that between a noun which is left unanchored and one which is brought into a more determinate relation with the discourse.

This may be seen from a simple contrast. In the sentence below the object is construed as a particular loaf or portion of bread already identifiable in the context:

\ex
\begingl
\gla Travat Marku piaštnik.//
\glb bread.\Dir{}.\Def{} Marek.\Trs{} eat-\Pf{}//
\glft \trsl{Marek ate the bread.}//
\endgl
\xe

If, however, the noun appears without the determiner, the reference is left less fixed and may be understood in a weaker or less individuated sense:

\ex
\begingl
\gla Trava Marku piaštnik.//
\glb bread.\Dir{} Marek.\Trs{} eat-\Pf{}//
\glft \trsl{Marek ate bread; Marek ate some bread.}//
\endgl
\xe

It must be observed that this suffix does not serve to indicate grammatical relation as such. It does not mark subjecthood, nor is it imposed by the mere fact that a noun stands as a core argument. In transitive clauses of the perfective or retrospective, it is indeed especially frequent with direct-case patient nouns; but this is because such nouns are often conceived as individuated and discourse-recoverable, not because the syntax requires the suffix. The same determiner may, where the sense calls for it, be found upon nouns in other cases also; yet outside the direct case such use is generally less neutral and more plainly expressive of topical, resumptive, contrastive, or otherwise salient reference.

The vocative excludes the determiner. In forms of direct address the person addressed is already sufficiently identified by the speech-act itself, and no further anchoring is needed. Generic expressions likewise dispense with the suffix as a rule, the bare noun being in such contexts the ordinary vehicle of the class or kind as such.

Beside this determiner the language possesses two preposed particles, \ird{si} and \ird{druh}, which are employed when the speaker wishes to introduce a referent as particularized, yet not as already established in the common ground. These particles thus stand, in point of meaning, between the wholly bare noun and the noun provided with the determiner. They do not make the noun definite; rather, they indicate what may be called specific indefiniteness.

Of these two particles, \ird{si} is the lighter and more usual. It answers approximately to \trsl{a certain} or \trsl{some particular}, and is the common means of introducing a referent which the speaker has in view, though the hearer has not yet been led to identify it.

\ex
\begingl
\gla si byl//
\glb \Spec{} child.\Dir{}//
\glft \trsl{a certain child; some particular child}//
\endgl
\xe

The particle \ird{druh} is weightier in tone and somewhat more selective in force. It often suggests \trsl{a given} or \trsl{one particular}, and may accordingly appear where the speaker means to single out the referent with greater insistence, or where the style is more elevated or deliberate.

\ex
\begingl
\gla druh byl//
\glb \Spec{} child.\Dir{}//
\glft \trsl{a certain particular child; one given child}//
\endgl
\xe

The difference between \ird{si} and \ird{druh} is not absolute, and their spheres overlap. In many passages either may stand without serious change of denotation. Still, the general tendency is plain enough: \ird{si} is slighter, more colourless, and more readily employed in ordinary discourse, whereas \ird{druh} is stronger, more selective, and not seldom somewhat more literary in effect.

These particles do not ordinarily co-occur with the determiner \ird{-ot/-t}. The reason is semantic. A noun marked by \ird{si} or \ird{druh} is presented as particularized indeed, but not yet anchored; a noun marked by \ird{-ot/-t} is presented as already anchored and recoverable. The two constructions are therefore, in the normal language, held apart.

Number does not directly alter this opposition. The noun itself is not obligatorily inflected for singular and plural. Plurality, where it is to be expressly marked, is indicated by the proclitic \ird{ne}, which stands at the left edge of the noun phrase. The plural marker and the determiner are independent of one another in form and function: the former denotes plurality, the latter anchoring.

\ex
\begingl
\gla ne byl \qquad ne bylot//
\glb \Pl{} child.\Dir{} \qquad \Pl{} child.\Def{}//
\glft \trsl{children \qquad the children}//
\endgl
\xe

The particles of specific indefiniteness may likewise accompany plural expressions:

\ex
\begingl
\gla ne si byl \qquad ne druh byl//
\glb \Pl{} \Spec{} child.\Dir{} \qquad \Pl{} \Spec{} child.\Dir{}//
\glft \trsl{some particular children \qquad a certain set of children}//
\endgl
\xe

When \ird{ne} is joined with a numeral or other quantifying word, it no longer denotes mere plurality, but rather an approximate number. In such combinations the determiner, as well as the particles \ird{si} and \ird{druh}, is ordinarily avoided, since approximate quantified expressions are not commonly treated as either anchored or narrowly individualized noun phrases.

A word must also be added concerning demonstratives. In the unmarked construction a demonstrative stands without the postposed determiner. Nevertheless, the language permits, in emphatic or contrastive style, the combination of demonstrative and determiner in the same noun phrase. Such double determination is not neutral, but serves to single out the referent with unusual force, whether resumptive, affective, or contrastive.

\ex
\begingl
\gla ša byl \qquad ša bylot//
\glb \Dem{}.\Prox{} child.\Dir{} \qquad \Dem{}.\Prox{} child.\Def{}//
\glft \trsl{this child \qquad this very child; this particular child here}//
\endgl
\xe

The full semantic range of the determiner, together with its relation to deixis, anaphora, topicality, contrast, and the introduction and subsequent tracking of referents, will be considered more fully in Chapter~\ref{ch:discourse}. The formal paradigm of the demonstratives themselves is given in Section~\ref{sec:demonstratives}.
\section{Number}

Iridian nouns are not formally inflected for number. A noun in the direct case form may be interpreted as singular or plural depending on context. The plural particle \ird{ne} may be used to explicitly mark plurality; it is a proclitic attaching to the leftmost element of the noun phrase.

\ex
\begingl
\gla ne byl//
\glb \Pl{} child//
\glft \trsl{children}//
\endgl
\xe

The use of \ird{ne} is optional. In contexts where number is recoverable from numerals, quantifiers, or shared knowledge, the particle is typically omitted. When used with a numeral or other quantifier, \ird{ne} signals approximation rather than plain plurality: \ird{ne hrona byl} means \trsl{about three children}, not \trsl{three children.}

\ird{Ne} is a proclitic and attaches to the front of the noun phrase, including any preceding demonstratives or modifiers. It cannot appear inside a possessive noun phrase between an oblique possessor and the head noun. \ird{Ne} cannot be combined with a \ird{zad-} (negative existential) clause.

\section{Uses of the indirect case}\label{sec:indirect-uses}

Beyond its role as the patient case of transitive clauses in the nominative-alignment aspects, the indirect case is the general non-core case of the noun phrase. It covers the functions that in older descriptions were distributed across several narrower cases. In modern analysis these uses form a single domain: the indirect marks dependence on another noun, relation to a source or whole, and a wide range of peripheral participant roles.

\subsection{Possession and nominal dependence}

The most basic non-core use of the indirect is to mark a dependent noun within the noun phrase, especially a possessor. The dependent noun precedes the head noun and forms a tight constituent with it.

\pex
\irdp{Marcí dum}{Marek's house}\\
\irdp{mámcí hašek}{my mother's bag}\\
\irdp{ša študencí tóm}{this student's book}
\xe

Demonstratives and other phrasal modifiers precede the entire possessive phrase and do not intervene between the oblique possessor and the head noun. The plural clitic \ird{ne}, by contrast, is placed immediately before the element whose plurality it marks.

\pex
\a \irdp{ša študencí tóm}{this student's book}
\a \irdp{ně študencí tóm}{the students' book}
\a \irdp{študencí ně tóm}{the student's books}
\xe

The same oblique dependency is used more broadly for close nominal relations that English often renders with \trsl{of}: kinship, authorship, association, and other kinds of inherent or identifying linkage.

\subsection{Material and partitive relations}

The indirect also marks the substance from which something is made and the noun that supplies the whole or source domain in a partitive relation. In this use the head noun denotes an object, measure, or portion, while the indirect noun names the material or larger whole from which it is drawn.

\ex
\irdp{kuní prosc}{silver spoon}
\xe

\ex
\begingl
\gla Ona pive štava unarížčem.//
\glb one beer-\Ind{} mug-\Dir{} \Refl{}-order-\mk{av-pv-1s}//
\glft \trsl{I ordered a mug of beer.}//
\endgl
\xe

Where the construction refers to an unbounded substance or to a portion abstracted from a whole, the indirect is the normal choice. By contrast, when the language counts discrete members of a set, Iridian prefers an overt quantificational construction rather than a bare oblique partitive dependent.

\pex
\a
\begingl
\gla Kroumaště na po zma jeca pivo.//
\glb refrigerator.\Ind{} in still few \Exst{}-\Cont{} beer.\Dir{}//
\glft \trsl{There's still some beer left in the refrigerator.}//
\endgl
\a
\begingl
\gla Študencí na hroná.//
\glb student-\Ind{} in three//
\glft \trsl{three of the students}//
\endgl
\xe

\subsection{Source and movement away}

With verbs of departure, removal, and separation, the indirect marks the source from which motion proceeds. This source reading contrasts with the direct case, which in comparable verbal environments may instead foreground the affected place as a more central argument.

\pex
\a
\begingl
\gla Dumí palžek.//
\glb house-\Ind{} leave-\Pf{}//
\glft \trsl{I left the house.}//
\endgl
\a
\begingl
\gla Dum palzinek.//
\glb house.\Dir{} leave-\Pf{}//
\glft \trsl{I left the \emph{house}.}//
\endgl
\xe

The same source-oriented oblique extends naturally to nouns followed by postpositions of separation, origin, exclusion, and replacement.

\section{Pronouns and demonstratives}\label{sec:pronouns}

Iridian distinguishes first and second person by means of a small set of personal pronouns. There are no independent personal pronouns of the third person. Reference to third persons is instead managed by the demonstrative system, together with a reduced set of anaphoric forms derived from it. The personal pronouns are inflected for the three ordinary cases, and each has beside it a possessive clitic, of which more will be said presently.

\begin{table}[ht]
	\footnotesize\sffamily
	\caption{Personal pronouns and possessive clitics.}\label{tab:personal-pronouns}
	\medskip
	\begin{tabularx}{0.75\textwidth}{L l l l l}
		\toprule
		              & {\sc direct} & {\sc transitive} & {\sc indirect} & {\sc clitic} \\
		\midrule
		1\textsc{sg}  & \ird{dá}   & \ird{dam}  & \ird{že}   & \ird{-(e)m}  \\
		2\textsc{sg}  & \ird{já}   & \ird{jí}   & \ird{jena} & \ird{-(e)š}  \\
		1\textsc{pl}  & \ird{mé}   & \ird{mně}  & \ird{mi}   & \ird{-(e)m}  \\
		2\textsc{pl}  & \ird{tová} & \ird{tvy}  & \ird{tě}   & \ird{-(e)st} \\
		\bottomrule
	\end{tabularx}
\end{table}

The indirect case of the personal pronoun serves, among its other functions, to express the possessor, precisely as the indirect case of an ordinary noun may do elsewhere in the language (Section~\ref{sec:indirect-uses}). Thus \irdp{že dum}{my house} is formally no different in structure from the corresponding nominal possessor construction. The possessor stands before the head noun in the usual prenominal position.

Alongside this analytic construction there exists a shorter and in many contexts commoner possessive expression by means of enclitic forms attached to the right edge of the noun. The first-person clitic appears as \ird{-m} after a vowel-final stem and as \ird{-em} after a consonant-final stem; the second-person singular behaves in like manner as \ird{-š} or \ird{-eš}, and the second-person plural as \ird{-st} or \ird{-est}. Hence \irdp{dumem}{my house; our house}. It will be observed that the first singular and the first plural have here fallen together in the same clitic form, \ird{-(e)m}; the distinction, where it matters, is left to the context or else restored by the use of the free pronoun.

The free indirect form and the clitic are not mutually exclusive. On the contrary, the language not seldom employs both at once where emphasis, contrast, or explicitness is desired, as in \irdp{že dumem}{my house}. Much the same liberty is found when the noun is at the same time accompanied by a demonstrative: \irdp{tuto že tom}{that book near you of mine}, \irdp{tuto tomem}{that book near you, mine}, and \irdp{tuto že tomem}{that book near you of mine} are all possible, though they differ in weight and emphasis.

Third-person reference belongs properly to the demonstrative system. Iridian possesses three demonstrative stems, distinguished according to the familiar threefold deictic opposition of nearness to the speaker, nearness to the addressee, and distance from both. These may be termed, in the usual fashion, proximal, medial, and distal.

\begin{table}[ht]
	\footnotesize\sffamily
	\caption{Demonstrative paradigm.}\label{tab:demonstratives}
	\medskip
	\begin{tabularx}{0.65\textwidth}{L l l l}
		\toprule
		           & {\sc direct} & {\sc transitive} & {\sc indirect} \\
		\midrule
		Proximal   & \ird{tak}  & \ird{tam}   & \ird{ten}  \\
		Medial     & \ird{tuto} & \ird{toho}  & \ird{taly} \\
		Distal     & \ird{je}   & \ird{jaky}  & \ird{jení} \\
		\bottomrule
	\end{tabularx}
\end{table}

In their primary and local use, these forms answer to the ordinary distinctions of deixis. Thus \irdp{je tom}{that book yonder} contains the distal form, whereas \irdp{tuto pect}{that cup near you} contains the medial. The indirect case of the demonstrative may further serve as a possessor in the same manner as the indirect of a noun or personal pronoun. One may therefore say \irdp{ten dum}{his or her house}, \irdp{taly pect}{his or her cup}, or \irdp{jení vrad}{his, her, or their town}. In such expressions the deictic opposition is preserved, but it is readily extended into discourse. The choice of demonstrative may indicate not only where the possessor stands in relation to speaker and hearer, but also how the referent is situated in the flow of the discourse.

This latter point is of some importance. When used pronominally, the three demonstratives do not merely point in space; they also arrange third-person referents according to accessibility and prominence. The proximal is naturally employed of the current topic or most salient participant; the medial may denote a secondary but still active referent, often one associated with the addressee’s sphere; and the distal tends toward what may conveniently be called obviative use, namely a referent less central, more displaced, backgrounded, or set over against another.

\begin{table}[ht]
	\footnotesize\sffamily
	\caption{Discourse function of the demonstrative series.}\label{tab:dem-discourse}
	\medskip
	\begin{tabularx}{0.9\textwidth}{L L L}
		\toprule
		{\sc form}   & {\sc spatial value}    & {\sc discourse value} \\
		\midrule
		Proximal & near speaker       & current topic, most accessible referent \\
		Medial   & near addressee     & addressee-linked or secondary active referent \\
		Distal   & away from both     & obviative, backgrounded, displaced, or contrastive referent \\
		\bottomrule
	\end{tabularx}
\end{table}

By this means Iridian is able to keep apart several third-person participants without the aid of verbal agreement. The forms are not mere substitutes for a single abstract \trsl{he} or \trsl{she}, but carry with them a definite discourse colouring.

\pex
\begingl
\gla Tak vrbnik, tuto spasnik, je těsnik.//
\glb \Prox{}.\Dir{} stand.up-\Pf{}, \Med{}.\Dir{} sit.down-\Pf{}, \Dist{}.\Dir{} leave-\Pf{}//
\glft \trsl{He stood up, he sat down, that one left.}//
\endgl
\xe

Alongside the full demonstratives there exists a reduced anaphoric set, phonologically derived from the distal series. These forms are used only pronominally, and only where the referent is already highly active and unmistakable in the discourse. They do not properly point; they merely resume.

\begin{table}[ht]
	\footnotesize\sffamily
	\caption{Reduced anaphoric forms (from distal demonstrative).}\label{tab:anaphoric-pronouns}
	\medskip
	\begin{tabularx}{0.6\textwidth}{L l l}
		\toprule
		{\sc case}   & {\sc full distal} & {\sc reduced} \\
		\midrule
		Direct     & \ird{je}   & \ird{i}  \\
		Transitive & \ird{jaky} & \ird{ji} \\
		Indirect   & \ird{jení} & \ird{ni} \\
		\bottomrule
	\end{tabularx}
\end{table}

The distinction is chiefly one of pragmatic force. The full demonstrative retains its deictic or contrastive value, and may be chosen whenever the speaker wishes to indicate distance, opposition, or a stronger gesture of reference. The reduced form, by contrast, is lighter and more nearly a pure anaphor. In broad terms one may say that \ird{tak} points, whereas \ird{i} merely carries forward an already established referent.

In ordinary connected speech, moreover, Iridian shows a marked tendency to omit pronouns wherever they are recoverable from the situation or the preceding discourse. It is, in this sense, a strongly pro-drop language. Overt personal pronouns are introduced chiefly for contrast, insistence, disambiguation, or some comparable rhetorical purpose. This tendency is especially notable in the second person. Direct pronominal reference to the addressee is relatively marked, and may suggest impatience, familiarity, or an undesirable directness. Courteous usage therefore prefers, wherever an overt expression is required, some noun of address instead: a personal name, a title, a kin term, the designation of an office, or the like. The most unmarked course is commonly silence; next in neutrality comes the respectful nominal expression; the overt second-person pronoun is the strongest and least neutral of the available choices.

There remains to be noticed the reflexive. Iridian possesses a distinct reflexive anaphor, whose forms show some remodelling under Slavic influence. This element is not a free pronoun of reference like the demonstratives, but a bound anaphor whose office is solely to mark local coreference.

\begin{table}[ht]
	\footnotesize\sffamily
	\caption{Reflexive anaphor forms.}\label{tab:reflexive}
	\medskip
	\begin{tabularx}{0.5\textwidth}{L l}
		\toprule
		{\sc case}   & {\sc form} \\
		\midrule
		Direct     & \ird{svem} \\
		Transitive & \ird{svě}  \\
		Indirect   & \ird{svy}  \\
		\bottomrule
	\end{tabularx}
\end{table}

The reflexive does not introduce a new discourse participant, nor does it rank one referent against another. It merely indicates that a given argument is coreferential with the local controller. For that reason it does not compete with the demonstratives in the ordinary tracking of discourse referents. The contrast is plain in the following examples:

\pex
\a
\begingl
\gla Svem Marku vednik.//
\glb \Refl{}.\Dir{} Marek.\Trs{} see-\Pf{}//
\glft \trsl{Marek saw himself.}//
\endgl
\a
\begingl
\gla I Marku vednik.//
\glb 3.\Dist{}.\Dir{} Marek.\Trs{} see-\Pf{}//
\glft \trsl{Marek saw him.} \lingcontext{non-reflexive; referent distinct from Marek}//
\endgl
\xe

The same opposition appears in possessive expressions. A demonstrative possessor remains non-reflexive; it refers to some other salient discourse participant, not back to the controller of the clause.

\pex
\a
\begingl
\gla Jení dum Marku vednik.//
\glb \Dist{}.\Ind{} house.\Dir{} Marek.\Trs{} see-\Pf{}//
\glft \trsl{Marek saw his house.} \lingcontext{someone else's house, not Marek's}//
\endgl
\xe

This distinction proves particularly useful where more than one third-person referent is simultaneously in play. The demonstrative system can separate discourse roles and degrees of accessibility; the reflexive, for its part, cuts through ambiguity by marking local identity directly.

\pex
\a
\begingl
\gla Je tam vednik.//
\glb \Dist{}.\Dir{} \Dist{}.\Trs{} see-\Pf{}//
\glft \trsl{He\textsubscript{1} saw him\textsubscript{2}.} \lingcontext{two distinct referents}//
\endgl
\a
\begingl
\gla Svem tam vednik.//
\glb \Refl{}.\Dir{} \Dist{}.\Trs{} see-\Pf{}//
\glft \trsl{He saw himself.} \lingcontext{reflexive; agent and patient coreferential}//
\endgl
\xe

\subsection{Interrogatives}

% [Port from 2020 archive — jede, ježe, jehát, jehu, jemí, jena, jak, zajehu, jiká, jišká, jeně, jení]

\subsection{Indefinite and negative pronouns}

% [Port from 2020 archive — nět/zide, niže/niho, etc.; note niho is superseded by zad- in the current analysis]

\section{Postpositions and case governance}\label{sec:prepositions}

Iridian is strictly head-final in its adpositional phrases. The governed noun phrase precedes the relational marker, and the marker itself is therefore best analyzed as a \term{postposed clitic adposition}. These items are phonologically weak, lean prosodically on the preceding noun phrase, and attach to the full governed phrase rather than to the noun stem alone. Orthographically they are here written as separate words for ease of reading, but morphologically they behave like enclitics.

All such clitic adpositions govern the \term{indirect case}. Older descriptions sometimes distributed these relations across several case labels or treated some of the markers as prepositions, but the synchronic generalization is simpler: the dependent noun phrase is inflected in the indirect, and the following clitic specifies the relation. The differences among the forms are therefore semantic, not differences in case government.

\subsection*{Morphological behavior}

The clitic attaches to the right edge of the governed noun phrase as a whole. A noun with modifiers, determiners, or possessors is therefore inflected first for the indirect case, after which the adposition follows as a phonological dependent. The syntax remains head-final throughout: the noun phrase comes first, and the clitic adposition closes the phrase.

\subsection*{Inventory}

The common clitic adpositions are \ird{na} \trsl{in, inside, on, at}, \ird{de} \trsl{into, to}, \ird{my} \trsl{for}, \ird{u} \trsl{near, by}, \ird{še} \trsl{with}, \ird{z} \trsl{from, out of}, \ird{nam} \trsl{without}, \ird{pale} \trsl{instead of}, and \ird{ahte} \trsl{except for}. Location, goal, proximity, benefit, accompaniment, source, exclusion, and replacement thus all belong to the same formal pattern: an indirect noun phrase followed by a phonologically weak clitic adposition.

\subsection*{Spatial and goal relations}

\pex
\a \irdp{dumě na}{in the house; at home}
\a \irdp{hradě de}{into the town; to the town}
\a \irdp{hradě u}{near the town}
\xe

The same pattern extends naturally to more specific spatial and path relations. In addition to the broad locative and goal clitics, the language also admits more specialized forms such as \ird{pod} \trsl{under, beneath}, \ird{nad} \trsl{over, above}, \ird{měz} \trsl{between, among}, \ird{do} \trsl{up to, as far as}, and \ird{mně} \trsl{through, via}. These too follow an indirect host noun phrase and behave as phonologically weak enclitics.

\pex
\a \irdp{stolě pod}{under the table}
\a \irdp{mostě nad}{over the bridge}
\a \irdp{domě měz}{between the houses; among the houses}
\a \irdp{hradě do}{up to the town}
\a \irdp{lesě mně}{through the forest; via the forest}
\xe

\subsection*{Participant and dependency relations}

\pex
\a \irdp{bylam my}{for the child}
\a \irdp{Jancí še}{with Janek}
\a \irdp{Marcí z}{from Marek}
\a \irdp{záhárí nam}{without sugar}
\a \irdp{Jancí pale}{instead of Janek}
\a \irdp{Jancí ahte}{except for Janek}
\xe

\subsection*{Clausal examples}

In full clauses the same structure is preserved: the host noun phrase is indirect, and the following adposition is glossed by its nearest English equivalent.

\pex
\a
\begingl
\gla Kroumaště na po zma jeca pivo.//
\glb refrigerator.\Ind{} in still few \Exst{}-\Cont{} beer.\Dir{}//
\glft \trsl{There's still some beer left in the refrigerator.}//
\endgl
\a
\begingl
\gla Bylam my Jancí še stóžách.//
\glb child-\Ind{} for Janek-\Ind{} with go-\Ctp{}//
\glft \trsl{(I am) going for the child with Janek.}//
\endgl
\a
\begingl
\gla Marcí z houba.//
\glb Marek-\Ind{} from gift.\Dir{}//
\glft \trsl{a gift from Marek}//
\endgl
\xe

The governing relation is thus constant across the system: the noun phrase takes the indirect, while the following clitic adposition contributes the relational meaning. This analysis captures both the head-final syntax of the phrase and the reduced phonological status of the adposition itself.

\section{Noun phrase structure}\label{sec:np-structure}

% ============================================================
% STATUS: STUB — to be developed
% The NP in Iridian is head-final: modifiers precede the head noun.
% Order within the NP:
%   ne (plural) > demonstrative > oblique possessor > stative-verb modifier > head
% Stative verbs as attributives: this is the primary adjectival strategy.
% The attributive suffix -ní can nominalize/adjectivalize a stative verb
% but the stative converb is also used.
% ============================================================

The canonical Iridian noun phrase is head-final. Modifiers of all kinds precede the head noun. The general ordering within the noun phrase is:

\medskip
\ird{ne} (plural) $>$ demonstrative $>$ oblique possessor $>$ property-root describer / stative-verb modifier $>$ head noun
\medskip

Possessors therefore occupy the same prenominal dependent slot described in Section~\ref{sec:indirect-uses}: they stand directly before the head noun and remain inside the noun phrase, while demonstratives and the plural clitic attach outside the possessive grouping.

Iridian has a very small closed set of lexical property roots used as simple
adnominal describers. These roots denote especially salient notions such as
\trsl{good}, \trsl{beautiful}, and \trsl{bad}, and they occur directly before
the head noun. Their use is lexically restricted and non-productive: speakers
do not freely coin new property-root describers for ordinary descriptive
content.

Outside this small residue, stative verbs (see Section~\ref{sec:stative} in
Chapter~\ref{ch:verbs}) serve as the primary attributive modification
strategy. Where English uses an adjective before a noun, Iridian normally uses
either a stative verb in an attributive/converbial form or a nominalized
stative verb. In many cases a property root and a stative verb form a lexical
pair: the root functions as an adnominal describer, while the corresponding
stative verb serves as the ordinary predicate.

% [Examples to be added]

\section{Event and participant nominals}\label{sec:event-nominals}

Iridian can form nouns directly from verbs. The resulting forms — called
\term{event nominals} and \term{participant nominals} — behave like
ordinary nouns in every grammatical respect, taking case endings, the
definiteness marker, and the plural clitic on exactly the same terms as
any other noun. What distinguishes them is their meaning: an event nominal
names the action or event itself (roughly what English does with
\trsl{the eating} or \trsl{the running}), while a participant nominal names
a person or thing defined by the role they play in that event (as English
does with \trsl{the eater} or \trsl{the food}).

Both types are formed by adding the suffix \ird{-ou} to a verb stem.
The details of which stem form to use in each case are treated in
Chapter~\ref{ch:verbs}; here the concern is with the resulting nouns
as noun phrases: how they inflect, how they interact with definiteness, and
how they encode the participants of the original event.

\subsection{Inflection and definiteness}

Nouns ending in \ird{-ou} follow the inflectional pattern of Class~II
nouns (see the section on inflectional classes above). In the direct case
the form ends in \ird{-ou}; the transitive ending is \ird{-ovy}; the
indirect is \ird{-ově}. The noun \ird{piaščikou} \trsl{the one who ate}
illustrates the full set:

\begin{table}[ht]
	\footnotesize\sffamily
	\caption{Inflection and definiteness for the participant nominal
	\ird{piaščikou}.}\label{tab:ou-inflection}
	\medskip
	\begin{tabularx}{0.72\textwidth}{L l l}
		\toprule
		{\sc case} & {\sc plain} & {\sc definite} \\
		\midrule
		Direct     & \ird{piaščikou}  & \ird{piaščikovat}  \\
		Transitive & \ird{piaščikovy} & \ird{piaščikovyt}  \\
		Indirect   & \ird{piaščikově} & \ird{piaščikovět}  \\
		\bottomrule
	\end{tabularx}
\end{table}

The definiteness marker (see Section~\ref{sec:definiteness}) attaches
after whichever case ending the noun already bears. The only irregularity
is in the direct case: the sequence \ird{-ou\,+\,-ot} does not surface as
\ird{*-out} but contracts to \ird{-ovat}. In the other cases the marker is
added without contraction: \ird{-ovy} yields \ird{-ovyt}, and \ird{-ově}
yields \ird{-ovět}.

\subsection{Two kinds of verb-derived noun}

Among nouns built on \ird{-ou}, it is useful to draw a line between two
kinds.

The first kind is the \term{lexical} or \term{resultant nominal}: a
noun tied to a specific verb that names the typical result, product, or
correlate of that verb's action. The verb for \trsl{to eat} can produce
\ird{piaštou} with the conventionalised meaning \trsl{food} — not any
particular act of eating, but the thing characteristically eaten or
produced by eating. Such forms must be learned verb by verb; many verbs
lack them entirely, relying instead on unrelated or inherited nouns.

The second kind is the \term{gerund}: the fully productive way of
turning any verb into a noun that names the event itself. Gerunds are
formed freely from any verb and always refer to the action as an abstract
event rather than to a fixed result. They correspond most closely to the
English nominalisation in \trsl{-ing} when used as a noun (\trsl{the
eating}, \trsl{the running}, \trsl{the remembering}).

Gerunds admit a further distinction based on what they contain. A
\term{simple} gerund names the event without specifying who was involved:
it means just \trsl{the eating}, nothing more. A \term{complex} gerund
names the event and also expresses the participants of that event inside the
noun phrase: \trsl{Marek's eating of the bread}, \trsl{Janek's
remembering}. No separate morphology marks the difference; the same
\ird{-ou} suffix produces both. What makes a given gerund simple or complex
is simply whether case-marked nouns appear inside it.

\subsection{Complex gerunds and their participants}

When a gerund expresses its participants internally, both the doer
(\textit{agent}) and the thing acted upon (\textit{theme}) appear in the
indirect case. This is the same case used throughout the noun phrase for
possession and dependent relations (see Section~\ref{sec:indirect-uses}).
The pattern is consistent with the rest of the grammar: just as
\ird{Marcí dum} means \trsl{Marek's house}, with the possessor Marek
in the indirect case, the doer of a gerundival event also appears in the
indirect case when expressed inside the noun phrase.

One consequence is that a complex gerund containing both a doer and a
theme will have two indirect-case nouns with no formal difference between
them. Consider:

\ex
\ird{Jancí podohletou} --- \trsl{the act of remembering Janek}
\textit{or} \trsl{Janek's act of remembering}
\xe

This is a genuine ambiguity. The indirect-case noun \ird{Jancí} (Janek)
could be either the person doing the remembering or the person being
remembered. Iridian does not resolve this ambiguity by grammatical rule.
A reader or listener is expected to use context, and this is by design:
complex event nominals regularly carry a degree of referential looseness
that a full finite clause would not. The grammar treats this as a natural
property of the construction rather than a defect.

When the speaker wishes to be explicit, two strategies are available.
The first is \term{word order}: the doer-indirect precedes the
theme-indirect. In a construction of the form \textit{A B gerund}, A is
the doer and B is the theme. The second is the \term{definiteness
marker}: a theme noun bearing the \ird{-ot/-t} suffix
(Section~\ref{sec:definiteness}) is read unambiguously as the theme and
not as the doer. Since the \ird{-ou} direct ending contracts with \ird{-ot}
to give \ird{-ovat}, this marker is visually distinctive when the theme
itself is a nominalized form.

\subsection{Participant nominals}

Participant nominals use the same \ird{-ou} suffix, but instead of naming
an event they name a \textit{participant} in that event, identified by the
role they occupy. The role a given noun names is determined by which verb
stem the suffix is attached to.

Iridian inflects verbs to foreground different participants as their
central argument (the system is described in Chapter~\ref{ch:verbs}).
Attaching \ird{-ou} to a given stem form yields a noun naming whichever
participant that form places at the centre. Three stem types are relevant
here:

\begin{itemize}

	\item The \term{antipassive stem} is the form of a transitive verb
	in which the doer of the action is treated as the central argument and
	the thing acted upon is expressed as a peripheral, indirect-case
	dependent. Adding \ird{-ou} to this stem produces an
	\term{agent-oriented} participant nominal meaning \trsl{the one who
	does X}. All five indicative aspect stems are available here. Two of
	them show special surface effects already noted in Chapter~\ref{ch:verbs},
	and two more contract regularly with \ird{-ou}:
	the perfective antipassive \ird{-ek} yields \ird{-ikou} (so
	\ird{piaščikou} means \trsl{the one who ate}), and the continuous
	antipassive \ird{-a} yields \ird{-ova} (so \ird{piaščova} means
	\trsl{the one who eats} or \trsl{the one who is eating}, whether
	habitually or at the present moment). The retrospective \ird{-aní}
	yields \ird{-anou}, and the progressive \ird{-ime} yields \ird{-imou};
	hence \ird{piaščanou} \trsl{the one who has eaten},
	\ird{piaščimou} \trsl{the one who is in the act of eating}, and
	\ird{piaščachou} \trsl{the one who is about to eat}.

	\item The \term{plain transitive stem} — the ordinary, unmarked form
	of a transitive verb — produces a \term{patient-oriented} participant
	nominal meaning \trsl{the thing that X is done to}. In Iridian's default
	treatment of transitive verbs (discussed in Chapter~\ref{ch:verbs}),
	it is the thing acted upon that occupies the central argument position
	of the construction. Attaching \ird{-ou} to this stem therefore names
	the recipient of the action. Here too all five aspects are available:
	\ird{piaštnikou} means \trsl{that which was eaten}, \ird{piaštnanou}
	\trsl{that which has been eaten}, \ird{piaštnova} \trsl{that which is
	habitually eaten}, \ird{piaštnimou} \trsl{that which is being eaten},
	and \ird{piaštnachou} \trsl{that which is about to be eaten}. More
	concretely, such forms may be interpreted as \trsl{food} understood as
	the participant affected by eating, with the chosen aspect supplying the
	expected temporal contour.

	\item The \term{applicative stem} is the form of a verb that promotes
	an indirect participant — typically a beneficiary, recipient, or other
	non-patient role — into the central argument position. Attaching \ird{-ou}
	to this stem produces a noun referring to that promoted participant.
	Thus \ird{piaštébikou}, built on the applicative stem of the verb for
	\trsl{to eat}, means \trsl{the one who was fed} or \trsl{the one for
	whom eating was done}.

\end{itemize}

The logic throughout is the same: whichever participant a given voice- and
aspect-marked verb form treats as its central argument, attaching \ird{-ou}
to that form yields a noun that names that participant. Antipassive-based
nominals name the doer; plain-transitive-based nominals name the thing
acted upon; applicative-based nominals name the participant promoted by the
applicative. Aspect contributes the expected interpretation of the event
profile: completed (\ird{piaščikou}), experientially or result-relevant
past (\ird{piaščanou}), ongoing or habitual (\ird{piaščova}),
currently in progress (\ird{piaščimou}), or intended / about to occur
(\ird{piaščachou}).

Participant nominals are particularly useful when the speaker wants to
identify someone by what they do or what is done to them, rather than by a
separate descriptive noun. They therefore tend to appear in constructions
that introduce or identify discourse participants, a topic taken up in
Chapter~\ref{ch:discourse}.

\subsection{The attributive form in \ird{-ovní}}

Participant nominals are Iridian's primary means of expressing what English
conveys with the relative pronouns \trsl{who}, \trsl{which}, and
\trsl{that}. The \ird{-ou} nominal stands alone as a freestanding noun; its
counterpart, the \term{attributive form} in \ird{-ovní}, is built from the
same verb stem and serves as a modifier directly preceding a head noun. The
construction that results is Iridian's standard equivalent of a restrictive
relative clause.

The two forms share a single organising principle: a nominalized verb phrase
denotes whichever participant is structurally \textit{privileged} within that
clause package. The suffix \ird{-ou} or \ird{-ovní} identifies the privileged
participant; every other core argument within the same nominalized domain
appears in the indirect case. The nominalized clause is its own enclosed
domain, so this rule operates within it independently of any surrounding
structure. When the noun phrase headed or modified by \ird{-ovní} is
embedded in a higher clause, it takes whatever case that outer clause
requires.

The paradigm below uses \ird{vrava} (Class~III, \trsl{letter}), \ird{pres}
(Class~II, \trsl{man}), and the verb \ird{obroh-} \trsl{to read}:

\pex
\a \ird{vravat presy obrohnik.} --- \trsl{The man read the letter.}
\a \ird{vravět obrohnikou} --- \trsl{the one who read the letter}
\a \ird{vravět obrohnikovní pres} --- \trsl{the man who read the letter}
\xe

In (\lastx b), the participant nominal picks out the reader as its referent.
The letter appears in the indirect case — \ird{vravě}, with the definiteness
suffix giving \ird{vravět} — as a non-privileged argument inside the
nominalized domain. In (\lastx c), the same stem takes the attributive suffix
\ird{-ovní} and immediately precedes the head noun \ird{pres}. The head noun
takes whatever case the surrounding clause requires; the \ird{-ovní} form
does not change.

\subsubsection*{Definiteness}

The attributive form \ird{-ovní} qualifies the head noun as something defined
by a verbal notion but does not itself encode definiteness. Definiteness,
where intended, is shown by the ordinary suffix \ird{-t}/\ird{-ot} on the
head noun, following case-inflection in the regular way; the \ird{-ovní}
modifier is unaffected.

\subsubsection*{Nested attributives}

Because each nominalized clause is its own domain, attributives may be
stacked recursively. The inner nominalized phrase becomes an argument within
the outer nominalized clause, appearing there in the indirect case as
required. The example below, introducing additionally \ird{zasd-} \trsl{to
write} and \ird{naz} \trsl{woman}, embeds one attributive inside another:

\pex
\a \ird{vravat naza zasdnik.} --- \trsl{The woman wrote the letter.}
\a \ird{nazam zasdnikovní vrava} --- \trsl{the letter that the woman wrote}
\a \ird{[nazam zasdnikovní vravět] obrohnikovní pres} ---
   \trsl{the man who read the letter that the woman wrote}
\xe

In (\lastx b), considered in isolation, the head noun \ird{vrava} appears in
the direct case: the letter is the privileged participant of the inner domain
(the thing written), and the woman appears as \ird{nazam} in the indirect.
In (\lastx c), this same inner phrase is embedded as a non-privileged
argument of the outer nominalized clause. Accordingly the inner head noun
shifts to the indirect: \ird{vrava} surfaces as \ird{vravět} (indirect with
definiteness suffix). The head noun of an \ird{-ovní} phrase always takes the
case required by the clause it is embedded in, not the case it would bear if
freestanding.

\subsubsection*{Practical limits on embedding}

The system is formally recursive, but in practice speakers rarely go beyond
two or three levels of embedding, and a single embedded attributive is by far
the most common. Deeper nesting produces noun phrases that are unwieldy to
produce and difficult to parse in speech. Speakers typically rephrase heavily
embedded content as a sequence of simpler clauses or a converb chain rather
than extending the nominalization further.
